\section{Introduction} \label{sec:Introduction}

The demand for handling long contexts is rapidly growing. While softmax attention~\cite{vaswani2017attention} has become the de facto solution for modeling various types of data, its computational cost grows quadratically with sequence length, motivating extensive research into more 
efficient long-context modeling.

Recently, Test-Time Training (TTT)~\cite{sun2024learning} has emerged as a promising approach for efficient sub-quadratic sequence modeling. TTT extends the concept of recurrent states in RNNs to a small, online-adapted sub-network.
The parameters of this sub-network also referred to as fast weight~\cite{schlag2021linear}, as they are rapidly adapted online via self-supervised objectives to memorize in-context information. 
Numerous recent studies~\cite{wang2025testtimeregressionunifyingframework,behrouz2024titans,behrouz2025itsconnectedjourneytesttime, karami2025lattice} have explored various online objectives, optimizers, and architectures for fast weight networks.

Despite these efforts, existing TTT methods struggle to scale effectively to long contexts, primarily due to extremely low hardware utilization in their TTT layers (often below 5\% peak FLOPS on modern GPUs). 
This inefficiency is because of the usage of
small mini-batch sizes, i.e. updating fast weights every token or every 16 to 64 tokens, which is conventionally assumed to be more effective for in-context learning.
Such small mini-batch results in poor parallelism and low compute intensity, and presents significant challenges for hardware-efficient 
implementation, especially when using large, nonlinear fast weights, making it difficult to achieve non-trivial (above 10\%) FLOPs utilization. 


In this paper, we adopt the opposite strategy and introduce Large Chunk Test-Time Training (\methodname). 
\methodname{} leverages extremely large chunk (from 2048 to 1M tokens) as the basic unit to update the fast weight.
Since the tokens within each large chunk are treated as an unordered set, we further integrate window attention
into \methodname{} to capture local dependencies within the chunk.
\methodname{} significantly enhances parallelism, leading to substantially improved GPU utilization (up to 70\% on NVIDIA A100s)  with just a few dozen lines of pure PyTorch code (see Appendix~\ref{sec:ttt_pseudocode}).
This efficiency enables the scaling of non-linear fast weights to enhance the memory capacity. And simple implementation allows easy integration of more effective test-time optimizers, such as Muon~\cite{jordan2024muon}.

Furthermore, \methodname's large-chunk design is also natural to model diverse N-dimensional data as we can align chunk-size with the internal structure of the data (e.g., grouping tokens within an image or consecutive video frames as a chunk).






We extensively  validate \methodname{} on three tasks spanning different modalities and data structures: 
\begin{itemize}
\item \em{Novel View Synthesis}.
Our model is capable of processing up to $128$ input images 
at a resolution of $960\!\times\!536$ leading to a maximum of $1$M tokens, and  outperforms 3D Gaussian Splatting~\cite{kerbl20233d} in terms of rendering quality under such input scale.
\item \em{Language Modeling}.
Our model achieves competitive performance compared to SoTA methods such as DeltaNet~\cite{yang2024parallelizing}, even though a chunk structure is not explicitly present in language data.
\item \em{Autoregressive Video Diffusion}. 
We adapt a 14-billion-parameter bidirectional video diffusion transformer into an autoregressive model by incorporating \methodname{} with sliding window attention. This adapted model generates consistent videos up to 56,000 visual tokens.
\end{itemize}


To summarize, our approach establishes an efficient, scalable, and highly performant framework for long 
sequence modeling across diverse modalities. 
By removing the dependency on low-level, hardware-specific implementations, \methodname{} enables broader exploration of the architectural design space. We believe this can democratize research in efficient long-context modeling and inspire the development of more novel and effective designs.  


